\section{Related Literature}\label{sec:lit}

Our analysis connects six bodies of work: theory and evidence on the trade-off between firm growth and profitability; the valuation of growth, its duration and the practitioner rules that summarize it; frontier estimation; the economics of AI and automation and its effects on firms; markups, rents and concentration; and recent work on AI-assisted research. For each, we keep the canonical statement and add the work of 2018 to 2026 that bears on our results; working papers and preprints are identified as such.

\subsection{Growth and profitability: theory and evidence}

The idea that growth is costly at the margin is old. \citet{penrose1959theory} argued that expansion is bounded by the managerial services available to plan and absorb it, so faster growth raises unit costs. \citet{uzawa1969time} formalized this ``Penrose effect'' as a concave relation between investment and the growth it buys, and the adjustment-cost models of \citet{lucas1967adjustment} and \citet{hayashi1982tobin} gave it a neoclassical form in which marginal $q$ governs investment and equals average $q$ under linear homogeneity. The closest antecedents of our model are \citet{marris1963model,marris1964economic}, \citet{solow1971some} and \citet{odagiri1981theory}, in which the firm's objective (value, growth, or a managerial preference over both) selects a point on a frontier trading growth against profitability or valuation. They remain the canonical treatments. Our feasible set is the same kind of object: a transformation curve between growth and margin, with each Rule-of-$X$ line acting as a price line, so the score-maximizing point is a tangency. \citet{mueller1969theory,mueller1972life} extended the managerial view to conglomerate mergers and to a life-cycle account in which growth opportunities decline as firms mature. On the financial side, \citet{higgins1977how} showed that sustainable growth is pinned down by margin, asset turnover, leverage and payout, a link that our capital-intensity correction and our restatement in terms of economic profit generalize.

The growth technology also has antecedents in marketing science. Growth bought by spending and limited by the unserved market resembles the sales-response model of \citet{vidale1957operations} and the goodwill model of \citet{nerlove1962optimal}, and our optimal-spend condition, under which spend equals the spend elasticity times the net score value of the growth it buys, is a \citet{dorfman1954optimal} condition. Modern firm-level evidence supplies the structure behind that technology. Exporters grow after entry by building a customer base through marketing and advertising, with flat markups and a customer base that responds little to past sales, rather than through low introductory prices \citep{fitzgerald2024how}; successful entrants build market share by reaching more customers, with no markup differences between new and established markets \citep{argente2025how}; and in a working paper, a firm's market share is tied mainly to its number of customers and its markup only to its sales per customer \citep{afrouzi2023concentration}. Growth bought with spend at roughly constant margins per customer, separate from the margin intercept, is what \autoref{ass:growth} posits. Empirically, the link between growth and profit, in either direction, is weak and varies across samples \citep{coad2007testing,davidsson2009growing,federico2015heterogeneous,lee2014relationship}, and firm-level effects explain more of the variance in both than industry effects do \citep{goddard2009sources}, which cautions against reading industry averages as frontier positions.

The closest modern relatives of our growth technology are models of customer capital. Reaching additional consumers in a market involves rising marginal penetration costs \citep{arkolakis2010market}; with search frictions the customer base is a slow-moving state variable that firms build with selling and marketing spend, which links that spend to investment, markups and Tobin's $q$ \citep{gourio2014customer}; customer relationships can be modeled as capital acquired by marketing \citep{drozd2012understanding}; and new plants stay small because of demand-side fundamentals: their customer base must be accumulated, and active accumulation driven by past production dominates passive accumulation \citep{foster2016slow}. In the $q$-theory of advertising, the shadow price of brand capital governs advertising and brand equity is a sizable share of firm value \citep{vitorino2014understanding,belo2022decomposing}, and in the optimal-control models of advertising the costate is the marginal value of a unit of sales or share \citep{sethi1973optimal,sethi1977dynamic,feichtinger1994dynamic}. Measured short- and long-term advertising elasticities of the sales level are small \citep{sethuraman2011advertising}, a bracket on diminishing returns rather than an estimate of our spend elasticity, which applies to excess growth. Our growth technology, its first-order condition for spend and the reading of the multiplier as the marginal value of a customer are therefore not new: the growth technology is a customer-capital structure in a normalization by revenue and market size, and the firm's weight on growth in \autoref{prop:dynamic} is the marginal $q$ of revenue of this literature. What we add is the reading of the optimum as the maximizer of a Rule-of-$X$ score at that weight, the contrast between that weight and an investor's screening weight, and the result that neither can be recovered from the cross-section of chosen points. That result is an instance of the identification problem that \citet{marschak1944random} raised for production functions: chosen points of heterogeneous firms trace a locus of optima, not any firm's technology.

\subsection{Valuing growth: duration and the practitioner rules}

{\sloppy In the dividend-discount model of \citet{gordon1959dividends,gordon1962investment} and the valuation framework of \citet{miller1961dividend}, growth adds value only when the return on new investment exceeds the cost of capital, so any linear rule that adds growth to margin embeds an implicit exchange rate between the two; our \autoref{prop:general} makes that rate explicit in terms of economic profit, and in steady state it restates the constant-growth valuation ratio of the managerial and dividend-discount models \citep{marris1963model,gordon1962investment} in the form of a score. \citet{schwartz2000rational} showed that high valuations of loss-making internet firms can be rationalized by uncertain, mean-reverting revenue growth, and \citet{mccarthy2017valuing} valued subscription businesses from disclosed customer metrics. The sensitivity of value to growth in these models is a duration effect in the sense of \citet{macaulay1938some}; our derivation of the screening weight as a duration connects to the implied equity duration of \citet{dechow2004implied}. Cash flows of longer duration earn lower expected returns: long-duration stocks earn about 1.1\% a month less than short-duration stocks \citep{weber2018cash}, and the value, profitability and investment premia reflect a premium on near-term cash flows \citep{gormsen2023duration}. If the discount rate falls with duration, the firm's own weight on growth is higher than our constant-rate calculation gives, which widens its gap with screening weights. Short-termism points the same way: firms that just meet analyst forecasts have lower R\&D growth and subsequent innovation, and in an estimated growth model short-termism slows growth \citep{terry2023macro}, as a firm acting on a screening weight below its marginal value of revenue would. The saturation term in our model resembles Bass-type diffusion \citep{bass1969new}.\par}

The Rule of 40 itself is a practitioner construct. \citet{feld2015rule} popularized it after hearing it from a late-stage investor: revenue growth plus profit margin should reach 40\%. Consulting studies find that it is rarely met and is associated with valuation premia. \citet{depeyrot2018hacking} report that only 16\% of 86 public software firms beat it in every year from 2013 to 2017 and that firms above the line traded at roughly double the EV/revenue multiple; \citet{roche2021saas} reach similar conclusions for more than 200 software firms over 2011 to 2021. Later work raised or reweighted the threshold. \citet{bosche2022sixty} asked whether 60 is the new 40; \citet{bessemer2024rulex} proposed the Rule of X and recommend multiplying growth by about 2 for private firms and 2 to 3 for public firms before adding free cash flow margin; and \citet{vista2026rethinking} describe a Rule of 60 as an emerging aspiration for AI-enabled software firms. In the other direction, \citet{lipman2026ai} argue that inference adds variable cost to software margins and that firms reinvesting in AI may need to accept a Rule of 30 for a time. The only peer-reviewed academic study of the rule we located is \citet{lee2024evaluating}, who evaluates it as an investment screen on 1,756 SaaS firms over 2003 to 2022 and proposes a ``Rule of 65.'' We found no modern peer-reviewed study of the growth-profitability frontier of software firms beyond it.

Scores also depend on which margin is used. \citet{bradshaw2002gaap} document the widening gap between GAAP and Street earnings, \citet{barth2012why} study why pro forma and Street earnings continued to exclude stock-based compensation after SFAS 123R, and \citet{mohanram2020stock} find that analysts who exclude it issue optimistically biased price targets. This motivates our explicit treatment of stock-based compensation in the margin measure. The split of selling, general and administrative expense into growth spend and maintenance has a market-based counterpart: exit prices imply the long-lived share of SG\&A and the depreciation of R\&D, giving intangible stocks about 15\% smaller than standard parameters \citep{ewens2025measuring}, and an intangible-adjusted value factor prices assets at least as well as the traditional one \citep{eisfeldt2022intangible}.

\subsection{Frontier estimation}

Our empirical frontier draws on two traditions. Stochastic frontier analysis, introduced by \citet{aigner1977formulation} and \citet{meeusen1977efficiency} and surveyed by \citet{kumbhakar2000stochastic}, separates noise from one-sided inefficiency under parametric assumptions. Deterministic envelopment, beginning with the data envelopment analysis of \citet{charnes1978measuring}, avoids functional-form assumptions but is sensitive to outliers. Partial frontiers address that sensitivity: the order-$m$ frontier of \citet{cazals2002nonparametric} and the conditional-quantile frontier of \citet{aragon2005nonparametric}, both treated in \citet{daraio2007advanced}, are robust to extreme observations. Upper-quantile regression \citep{koenker1978regression} offers a parametric counterpart. We have not found prior applications of these estimators to the growth-margin plane of software firms.

\subsection{AI, automation and firms}

Our reduced-form mapping of AI effects is informed by task-based models of automation, in which automation displaces labor from tasks while new tasks reinstate it \citep{acemoglu2018race}; we do not model reinstatement. \citet{aghion2019artificial} show that automation can coexist with balanced growth because, as in \citet{baumol1967macroeconomics}, sectors that resist automation come to dominate expenditure and constrain aggregate growth. \citet{trammell2023economic}, \citet{davidson2021could} and \citet{erdil2023explosive} examine conditions under which AI could produce much faster growth, \citet{epoch2025gate} build an integrated assessment model of AI automation, and \citet{korinek2024scenarios} set out scenarios for the path to general AI, the precedent for presenting our horizons as scenarios. \citet{acemoglu2025simple} argues from task exposure and cost-savings estimates that aggregate productivity effects over the next decade are likely to be modest. Technical exposure also overstates what is worth automating: in a working paper on computer vision, only 23\% of the wages paid for exposed tasks would be economical to automate at current costs \citep{svanberg2024beyond}.

Firm-level evidence bears on the channels in Table~\ref{tab:ai-map}. \citet{brynjolfsson2025generative} document productivity gains from a generative AI assistant for customer-support agents, in an experiment access to a generative AI tool cut the time professionals took on writing tasks by about 40\% and raised quality \citep{noy2023experimental}, and \citet{eloundou2024gpts} measure occupational exposure to large language models. Gains are uneven: inside the frontier of AI capability, knowledge workers completed 12.2\% more tasks 25.1\% faster, while on a task outside it AI lowered performance \citep{dellacqua2026navigating}, and in a randomized trial early-2025 AI tools increased the time experienced developers took to complete tasks \citep{becker2025measuring}, a preprint. In labor markets, Danish administrative data show precise null effects on earnings and hours two years after the launch of ChatGPT \citep{humlum2025still}, while employment of workers aged 22 to 25 in the most exposed occupations fell by 13\% relative to others \citep{brynjolfsson2025canaries}; both are working papers. Digitization lowers the cost of verification among other costs \citep{goldfarb2019digital}; our accountability floor is the part of that cost that does not fall, the cost of verifying AI-run controls. Firms that invest in AI grow faster in sales, employment and value, mainly through product innovation, and that growth is concentrated among larger firms and associated with higher industry concentration \citep{babina2024artificial}, which supports the growth channels ($\lambda$ and $G$) alongside the cost channel ($a$) and cautions against reading free entry as the norm. Markets capitalized labor substitution at once: firms whose workforces were most exposed to generative AI outperformed by 5\% in the two weeks after ChatGPT's release \citep{eisfeldt2023generative}, and within firms an occupation's employment share falls with the mean exposure of its tasks (about 14.5\% over five years per standard deviation) and rises with the concentration of exposure (about 7.5\%), while firms that adopt AI, instrumented by a university-hiring shift-share, grow faster in revenue, productivity and employment \citep{hampole2025artificial}; both are working papers. Adoption inside firms is slow. Generative AI assisted only 1 to 5\% of U.S.\ work hours in late 2024 \citep{bick2026rapid}; firm use of AI in producing goods or services rose from 3.7\% of firms in September 2023 to 5.4\% in February 2024 \citep{bonney2024tracking}; the Census Bureau then broadened the question to use in any business function, which raised measured adoption by 47\% to 159\% in each of five industries \citep{allen2026monitoring,census2025btoswording}, and under the new question 18\% of firms (32\% weighted by employment) used AI in November 2025 to January 2026, 57\% of these adopters in three or fewer functions, with firm performance positively associated with the breadth of integration \citep{bonney2026microstructure}. Before generative AI, fewer than 6\% of U.S.\ firms had adopted the AI technologies of 2018 \citep{mcelheran2024ai}. This is the slow approach that our attainable frontier formalizes. Adoption comes in bundles rather than task by task when tasks are complements \citep{agrawal2024artificial,gans2026oring}, which is the logic of programs and waves in that frontier; optimal adoption of a technology whose risks are learned over time is gradual \citep{acemoglu2024regulating}, and a preprint obtains staged commitment from pilots that build firm-specific capability \citep{tewari2026pilot}. On costs, the price of reaching fixed benchmark performance fell between roughly 9 and 900 times per year depending on the task \citep{epoch2025inference}, and by about 5 to 10 times per year for models on the cost-performance frontier \citep{gundlach2025price}, while the cost of running the best-performing model on a benchmark rose 3 to 18 times per year as token volumes grew \citep{gundlach2025price}. The estimate we cite for the price elasticity of demand for inference, a preliminary $-1.1$ across providers of the same open model, measures substitution between providers, and its authors report that aggregate elasticities are substantially smaller \citep{demirer2025emerging}; whether spending on inference, and the compute share of our ceiling, rise or fall is therefore a scenario. Post-quantum cryptography adds a dated cost: NIST deprecates quantum-vulnerable public-key algorithms after 2030 and disallows them after 2035 \citep{moody2026transition}, and the only quantified migration cost is a rough order of magnitude for U.S.\ federal civilian systems \citep{omb2024report}.

\subsection{Markups, rents and concentration}

Work on markups and intangibles explains why high scores are possible at all: \citet{deloecker2020rise} document rising markups, \citet{haskel2017capitalism} emphasize the scalability of intangible assets, \citet{crouzet2022economics} and \citet{corrado2022intangible} set out how non-rival, partly excludable intangibles scale and how to measure them, and \citet{peters2017intangible} and \citet{crouzet2023rents} show how intangible capital and rents alter the relation between $q$ and investment. Our free-entry result formalizes the implication that, where entry is free, AI-driven cost reductions raise scores durably only if they accrue to scarce inputs. Its equilibrium is the homogeneous-product case of free entry with fixed costs and an integer number of entrants \citep{mankiw1986free}; with differentiated products the benchmark would be monopolistic competition \citep{dixit1977monopolistic}, and switching costs and network effects are one source of the non-replicable input that sustains a rent \citep{farrell2007coordination}. Rising markups have gone with the rise of superstar firms \citep{autor2020fall}. Which inputs stay scarce is itself a question of modern theory. Data-driven markets tend to tip when more user data lower the cost of quality \citep{prufer2021competing}, and the largest firms are valued mainly for their data \citep{farboodi2023data}; a data rent is durable when the market tips and fades when it does not \citep[see also][]{hagiu2023data}. Because data are nonrival, such a rent also requires that the data be excluded from rivals \citep{jones2020nonrivalry}, and competition in markets for training data, input data and predictions determines how much of it a firm keeps \citep{gans2026market}. At the level of the AI supply chain, the amortized cost of training the most compute-intensive models has grown about 2.4 times a year since 2016 \citep{cottier2024rising}, a preprint: a rising fixed cost of the kind that concentrates supply, as do economies of scale and scope in foundation models \citep{korinek2025concentrating}, and with it the exposure to a single provider that our common-mode failure term prices. That term has a long history in reliability engineering, where common-cause failures are the dependent failures that redundancy does not remove \citep{mosleh1989nureg4780,mosleh1998nureg5485} and the beta-factor model is the most common form \citep{hokstad2008common}; in software, versions developed independently still fail together \citep{eckhardt1985theoretical,littlewood1989conceptual,knight1986experimental}. Our accountability floor applies the term to stacked reviewers of AI output: layers built the same way do not deliver independence.

\subsection{AI-assisted research}

Automated research systems are advancing quickly \citep{lu2024ai,si2025can,gottweis2026accelerating}, and \citet{korinek2023generative} catalogs uses of generative AI in economics. arXiv requires authors to report significant use of generative AI, states that such tools should not be listed as authors, and holds authors responsible for all content \citep{arxiv2026policy}; Appendix~\ref{app:protocol} documents our compliance.

\subsection{Our contribution}

Relative to this literature, the growth technology and its first-order condition are standard. We read the Rule of 40 and the Rule of $X$ as price lines on a frontier of attainable growth and margin, distinguish the firm's own weight on growth (the marginal value of revenue) from the investor's screening weight and show that the cross-section of chosen points identifies neither, derive a score budget that bounds how long extreme scores can last, charge the full costs of running AI against the ceiling, including an accountability floor whose common-mode part does not vanish under automation, and define the frontier attainable by a given date when transitions are risky and capacity is limited. Empirically, we estimate the theory's quantile frontier on a panel of software firms with every item as first reported, under four definitions of the sample, and state a claim only where it holds in all of them. Unlike \citet{lee2024evaluating}, who evaluates the rule as an investment screen, we model the frontier it is applied to, and its limits.
